\section{The dictionary for posets}

\subsection{Alexandrov sheaves and hypercompletion}

\begin{statement}{A.1}{Known}{The exact Alexandrov comparison}
For every small poset $P$, right Kan extension from principal opens gives
an equivalence
\[
 R_P:\Fun(P^{\op},\Spaces)\xrightarrow{\ \simeq\ }
 \Sh(P_{\downarrow})^{\wedge},\qquad
 R_P(F)(U)=\limop_{p\in U^{\op}}F(p).\tag{A.1}
\]
The superscript $\wedge$ denotes hypercompletion. The composite into
$\Sh(P_{\downarrow})$ is fully faithful. It is an equivalence without
hypercompletion if $P$ is finite, or more generally if every principal
downset $P_{\leq p}$ is finite. For set-valued sheaves the comparison is an
equivalence for every $P$.
\end{statement}
\begin{proof}
The functor $p\mapsto U_p$ from $P$ to $\operatorname{Open}(P_{\downarrow})$
is fully faithful. Its induced Grothendieck topology is trivial: a covering
family of $U_p$ contains an open containing $p$, and any such open contains
all of $U_p$. The basis theorem for hypercomplete sheaves
\citep[Proposition 3.12.11 and Example 3.12.15]{BGH2020} consequently
identifies hypersheaves with arbitrary diagrams on these principal opens.
Right Kan extension has the displayed formula. Since $p$ is initial in
$(P_{\leq p})^{\op}$, its value on $U_p$ is $F(p)$; hence restriction is
an inverse. Presheaf $\infty$-topoi are hypercomplete, since limits and
truncations are objectwise. This also proves full faithfulness into
ordinary sheaves.

A finite Alexandrov space has a hypercomplete sheaf $\infty$-topos
\citep[Example 3.12.15]{BGH2020}. If each $U_p$ is finite, restriction to
this open cover detects equivalences and hypercompleteness, giving the
stated extension. The ordinary comparison theorem for a basis gives the
set-valued assertion without a finiteness assumption. The empty poset
gives the terminal category on both sides; disjoint unions give products.
\end{proof}

\begin{statement}{A.2}{Refuted}{Ordinary sheaves for arbitrary posets}
The removal of hypercompletion from A.1 for arbitrary posets is false.
It is still false if ``locally finite poset'' means only that every closed
order interval is finite.
\end{statement}
\begin{proof}
Let $Q=\mathbb N\times\{0,1\}$, with $(i,a)<(j,b)$ exactly when $i<j$,
and put $P=Q^{\op}$. Every interval is finite. Every upper open of $Q$ is
generated by at most two minimal elements, so finite gluing shows that
filtered colimits in $\Sh(Q_{\uparrow})$ are computed sectionwise.
Its terminal object is therefore compact.

Set $Q_n=\{0,\ldots,n\}\times\{0,1\}$ and
$F_n=\operatorname{Lan}_{Q_n\hookrightarrow Q}(*)$. Pointwise,
$\colim_n F_n=*$, whereas, at rank greater than $n$,
\[
 F_n(j,b)=BQ_n=(S^0)^{*(n+1)}\simeq S^n.
\]
If $*$ were compact in $\Fun(Q,\Spaces)$, it would be a retract of some
$F_n$. The essential image of this fully faithful left Kan extension is
closed under retracts, by its counit criterion. Thus $*$ would itself be
that left Kan extension, contradicting $S^n\not\simeq *$.
The terminal objects have different compactness properties, so the
categories cannot be equivalent. This is the counterexample of
\citet[Example A.13]{Aoki2023}, with the variance reversed.
\end{proof}

\begin{statement}{A.3}{Proved}{Presentable coefficients on the Alexandrov model}
Let $\cC$ be a presentable $\infty$-category. Define
$\Sh^{\wedge}(P_{\downarrow};\cC)$ by hyperdescent, equivalently by
testing hyperdescent after every functor $\Map_{\cC}(c,-)$.
Then
\[
 \Sh^{\wedge}(P_{\downarrow};\cC)\simeq\Fun(P^{\op},\cC),
\]
with the same limit formula as A.1. If every principal downset is finite,
ordinary $\cC$-valued sheaves give the same category.
\end{statement}
\begin{proof}
The presheaf Yoneda embedding $\cC\to\Fun(\cC^{\op},\Spaces)$ preserves
and detects limits. Apply the basis comparison of A.1 after
$\Map_{\cC}(c,-)$, for every $c$, to obtain the right Kan extension
formula and its inverse. The diagrams so obtained are represented in
$\cC$ because the requisite small limits exist. Hyperdescent and the
unit and counit equivalences are detected by these mapping-space
functors. Under the finiteness hypothesis, a space-valued sheaf on each
principal open is already hypercomplete; applying this to all these
mapping-space functors proves the last assertion. This argument uses
no compactness assumption on $c$.
\end{proof}

\subsection{Realization and the geometric qualification}

\begin{statement}{A.4}{Proved}{The minimum-vertex stratification}
For a finite poset $P$ let $Z_P=|NP|$ and let
\[
 s_P:Z_P\longrightarrow P_{\downarrow},\qquad
 s_P([p_0<\cdots<p_k;t_0,\ldots,t_k])
       =\min\{p_i:t_i>0\}.
\]
This is a continuous stratification. Each stratum is contractible,
and its exit-path $\infty$-category is equivalent to $P^{\op}$.
Consequently
\[
 \Sh_{P\text{-cons}}(Z_P;\Spaces)
 \simeq\Fun(P^{\op},\Spaces)
 \simeq\Sh(P_{\downarrow};\Spaces).
\tag{A.4}
\]
The statement also holds locally when the order complex is locally finite.
Here strata are allowed to be locally contractible polyhedra; they are
not required to be manifolds.
\end{statement}
\begin{proof}
For a downset $U$, the inverse image $s_P^{-1}U$ consists of points with
positive barycentric coordinate at some vertex in $U$, and is a union of
open vertex stars. The stratum of $p$ is the part of the cone
$\{p\}*|N(P_{>p})|$ on which the coordinate of $p$ is positive.
Moving this coordinate to one contracts the stratum to $p$.

We also verify conicality of this particular coarsening. Write
$C_\infty T$ for the open cone with radial coordinate in $[0,\infty)$,
including its apex when $T=\varnothing$. On the open vertex star
$\operatorname{St}^{\circ}(p)=\{t_p>0\}$, divide all other barycentric
coordinates by $t_p$. The coordinates below and above $p$ vary
independently, since a chain below $p$ and a chain above $p$ concatenate
with $p$ to a chain. This gives a homeomorphism
\[
 \operatorname{St}^{\circ}(p)
 \cong C_\infty|N(P_{<p})|\times C_\infty|N(P_{>p})|.
\tag{A.4a}
\]
Explicitly, the two radial coordinates are
$r_-=\sum_{q<p}t_q/t_p$ and $r_+=\sum_{q>p}t_q/t_p$; their angular
coordinates are the respective normalized barycentric coordinates.
The inverse has $t_p=(1+r_-+r_+)^{-1}$ and rescales each angular
coordinate by $r_\pm/(1+r_-+r_+)$. These formulas respect all face
identifications, and are continuous in the cone topologies. The stratum
of $p$ is the first apex times the second cone. Away from the first
apex the stratum is determined by the minimum in the lower link,
independently of the second factor. Thus (A.4a), with the factors
interchanged, is a conical chart over $P^{\op}$ in the usual upward-open
convention. The stratum factor is a locally contractible polyhedron.
Consequently the raw exit simplicial set is an $\infty$-category
\citep[Theorem A.6.4]{LurieHA}, and agrees with the intrinsic exit
category of this exodromic stratification
\citep[Recollection 5.6.1]{HPT2026}.

We give the categorical comparison explicitly. Let $\mathsf{Ch}(P)$ be
the face poset of nonempty finite chains, with arrows given by inclusion.
The cell stratification of $Z_P$ has exit category $\mathsf{Ch}(P)$:
an exit from an open face to an incident coface has a contractible space
of choices in a linear simplex, and the simplexwise choices and their
compositions glue. This is the simplicial exit calculation
\citep[Theorem A.6.10]{LurieHA}. Passing to the coarser stratification
localizes this category at inclusions preserving the minimum. The
localization functor is
\[
 m:\mathsf{Ch}(P)\to P^{\op},\qquad \sigma\mapsto\min\sigma.
\]
Apply the universally localizing last-vertex map to $N(P^{\op})$
\citep[Proposition 6.3.6.2]{KerodonSubdivision}. Its subdivision is
$N\mathsf{Ch}(P)$, and its last vertex is the minimum in $P$.
This proves the asserted localization including its mapping spaces.
Explicitly, an arrow $p\to q$ in $P^{\op}$, with $q\leq p$, is represented by
$\{p\}\to\{q,p\}\leftarrow\{q\}$, where the backward arrow is inverted.
For $r\leq q\leq p$, the chain $\{r,q,p\}$ supplies the composition
relation.
Every chain is identified with its minimum by
$\{\min\sigma\}\to\sigma$; the universal localization theorem supplies
the full coherence assertion beyond these elementary roofs.

The coarsening theorem for exodromy identifies this localization with the
intrinsic exit category of $s_P$
\citep[Theorem 3.3.5 and Theorem 5.1.7]{HPT2026}; the conical charts just
proved identify it with the raw exit category as well. The finite
polyhedron is paracompact and locally of singular shape, and the finite
stratifying poset satisfies the ascending chain condition. Hence
\citet[Proposition A.5.9 and Theorem A.9.3]{LurieHA} give (A.4) for
ordinary constructible sheaves and show that these sheaves are
hypercomplete.
For a locally finite order complex the argument applies on finite star
neighbourhoods and glues by open descent. No assertion about an arbitrary
infinite-dimensional realization is hidden in this last local argument.
\end{proof}

\begin{statement}{A.5}{Known}{The coefficient range of geometric exodromy}
For the finite geometric model in A.4, (A.4) holds with $\Spaces$ replaced
by any compactly assembled presentable $\infty$-category $\cC$, and also
by any stable presentable $\infty$-category. More generally, for every
$\cC$ there is a fully faithful functor
\[
 \Fun(P^{\op},\cC)
 \simeq \Sh_{P\text{-cons}}(Z_P;\Spaces)\otimes\cC
 \hookrightarrow \Sh_{P\text{-cons}}(Z_P;\cC).
\tag{A.5}
\]
For an arbitrary presentable $\cC$, its essential surjectivity is an
additional assertion, not a consequence of presentability alone.
\end{statement}
\begin{proof}
The inclusion of constructible space-valued sheaves in all sheaves is
both a left and a right adjoint for this exodromic stratification.
Tensoring with $\cC$ preserves its fully faithful adjunction and places
the image in the constructible objects. The identification of the tensor
product with the functor category follows from
\citet[Proposition 4.8.1.17]{LurieHA}.
The essential-surjectivity assertions are precisely
\citet[Lemma 4.1.2 and Corollary 4.1.15]{HPT2026}: compactly assembled
coefficients work in general, while stable presentable coefficients work
for a noetherian stratifying poset. Our poset is finite.
\end{proof}

\begin{statement}{A.6}{Open}{The unrestricted geometric clauses}
This audit does not establish that (A.5) is essentially surjective for
every presentable $\cC$, even for every finite poset $P$. Nor does it
establish, for an arbitrary small poset with its unrestricted weak CW
realization, that ordinary constructible sheaves on $Z_P$ are precisely
$\Fun(P^{\op},\Spaces)$. These are separate statements from the fully
proved Alexandrov hypersheaf assertion A.1.
\end{statement}
The attempted tensor argument stops at the fully faithful functor (A.5):
tensoring does not by itself identify all objects detected as locally
constant after restriction to strata. For the unrestricted realization,
the finite-star proof stops when such a star is no longer finite; ordinary
and hyperconstructibility can differ. The first coefficient instance not
settled by that argument is the two-element poset with a presentable
coefficient category outside the coefficient classes of A.5. The explicit
infinite poset in A.2 is the first realization test retained here; A.2
refutes its ordinary Alexandrov clause, not by itself its geometric clause.
A.12 nevertheless computes the raw exit category for every small poset;
the unproved assertion here concerns its classification of ordinary
constructible sheaves.
We have not constructed a counterexample to either assertion in this
paragraph.

\subsection{Restriction, extensions, and local systems}

\begin{statement}{A.7}{Proved}{The three functors and the comma formulas}
For a monotone map $u:P\to Q$ and presentable $\cC$, restriction has both
adjoints,
\[
 u_!\dashv u^*\dashv u_*:
 \Fun(P^{\op},\cC)\rightleftarrows\Fun(Q^{\op},\cC),
\]
where the displayed arrows denote the three appropriately directed
functors and
\[
 (u_!F)(q)=\colim_{(q\backslash P)^{\op}}F,
 \qquad
 (u_*F)(q)=\limop_{(P/q)^{\op}}F.
\tag{A.7}
\]
They are the corresponding adjoints on Alexandrov hypersheaves. On the
finite geometric models, in the coefficient ranges of A.5, restriction
is geometric inverse image. For arbitrary $\cC$ the same compatibility
holds on the fully faithful diagram images of (A.5). The other two
functors are the intrinsic adjoints on these diagram categories and,
when A.5 is an equivalence, on the whole constructible categories.
They are not, without further hypotheses, the ambient topological
pushforwards.
\end{statement}
\begin{proof}
These are the pointwise formulas for left and right Kan extension along
$u^{\op}$. An arrow $u^{\op}(p)\to q$ in $Q^{\op}$ is an arrow
$q\to u(p)$ in $Q$, giving the first index; the second is dual. For
hypersheaves the functors are transported through A.3. For finite
Alexandrov spaces, inverse image takes stalk $q$ to stalk $u(p)$, and
direct image has sections over $u^{-1}(U_q)=\{p:u(p)\leq q\}$, giving
the second formula directly. In the geometric model
$s_Q|Nu|=u s_P$; the functor on exit categories is $u^{\op}$, so exodromy
identifies inverse image with restriction. The uniqueness of adjoints
then identifies its intrinsic adjoints, but says nothing about preservation
of constructibility by ambient operations. Part D gives a finite
counterexample to that additional claim.

For $P=\varnothing$, the first formula gives an initial object and the
second a terminal object at every $q$. These formulas remain valid when
$u$ misses connected components. The symbol $u_!$ here means left Kan
extension, not a prior compact-support construction on topological spaces.
\end{proof}

\begin{statement}{A.8}{Proved}{Strict fibres versus comma categories}
For $q\in Q$, the strict fibre $P_q=\{p:u(p)=q\}$ has canonical full
inclusions
\[
 k_q:P_q\hookrightarrow q\backslash P,\qquad
 l_q:P_q\hookrightarrow P/q.
\]
The comparison
$\colim_{P_q^{\op}}F\to(u_!F)(q)$ is an equivalence for every diagram
on the comma category whenever $k_q^{\op}$ is homotopy final. The
comparison $(u_*F)(q)\to\limop_{P_q^{\op}}F$ is an equivalence for
every diagram whenever $l_q^{\op}$ is homotopy initial.
The precise tests are, respectively, weak contractibility of
\[
 \{x\leq p:u(x)=q\}\quad(q\leq u(p)),\qquad
 \{x\geq p:u(x)=q\}\quad(u(p)\leq q).
\tag{A.8}
\]
\end{statement}
\begin{proof}
For a point $(p,q\leq u(p))$ of $q\backslash P$, the comma category
testing finality of $k_q^{\op}$ has objects $x\in P_q$ with $x\leq p$.
Its nerve is the nerve of the first displayed poset, possibly with
opposition, which does not change weak contractibility. For limits the
initiality test gives the second poset in the same way. The homotopy
cofinality theorem now gives the comparisons, including their directions.
Empty comma categories fail the test. Testing only the total homotopy
type of $k_q$ or $l_q$ would be insufficient for all coefficients; the
pointwise comma tests are what is used. Part D identifies them with
smoothness and properness under the manuscript's variance.
\end{proof}

\begin{statement}{A.9}{Known}{Local systems and the reflection}
Let $S$ be any small category and let
$\ell:NS^{\op}\to BS$ be groupoid completion, using reversal to
identify the completions of $S$ and $S^{\op}$. For every presentable
$\cC$, restriction is fully faithful and has both adjoints:
\[
 \ell_!\dashv \ell^*\dashv\ell_*.
\]
Its image is exactly the diagrams taking every arrow to an equivalence.
Thus
\[
 \mathcal P(S)^{\lc}\simeq\Fun(BS,\Spaces),\qquad
 \HOT^{\lc}(S)\simeq\Ho\Fun(BS,\Spaces).
\tag{A.9}
\]
The reflection of VI ms.~65--69 is, by its universal property,
$\rho=\ell_!$ with inclusion $\ell^*$; its endofunctor on all diagrams
is $\ell^*\ell_!$. The right projector is $\ell^*\ell_*$.
\end{statement}
\begin{proof}
The universal property of $\ell$ identifies functors on $BS$ with exactly
the functors inverting all arrows of $NS^{\op}$, including their mapping
spaces. Small limits and colimits in $\cC$ give both Kan extensions.
The counit $\ell_!\ell^*\to1$ and the unit $1\to\ell_*\ell^*$ are
equivalences because $\ell^*$ is fully faithful. The manuscript's
reflector is characterized by the same adjunction and is therefore
uniquely equivalent as an adjoint, not literally equal as a representative
construction. Cisinski proves the corresponding assertion for arbitrary
derivators \citep[Theorem 3.13]{Cisinski2009}; the attached GALOIS A.6 and
A.8 state these comparisons with the same variance.

For the downset Alexandrov model this is local \emph{hyper}constancy.
In the finite setting it is ordinary local constancy: constancy on
$U_q$ forces $F(q)\to F(p)$ to be an equivalence for every $p\leq q$,
and conversely all such equivalences identify the restriction with a
constant diagram on $U_q$. The same interpretation on $Z_P$ follows
from (A.4). We do not identify local hyperconstancy with ordinary local
constancy on every infinite-dimensional space.
\end{proof}

\begin{statement}{A.10}{Proved}{A computation and the need for homotopy colimits}
For $P=[1]$, write a presheaf as $F=(B\xrightarrow{v}A)$. Its reflection
and coreflection onto local systems are the constant arrows on $A$ and
$B$, respectively. An ordinary colimit cannot replace a homotopy colimit
in the preceding formulas without an additional hypothesis on the diagram
or its indexing category.
\end{statement}
\begin{proof}
The category $[1]^{\op}$ has terminal object $0$ and initial object $1$,
so its colimit is $A$ and its limit is $B$. The unit of the reflection
has components $1_A$ and $v$; the counit of the coreflection has components
$v$ and $1_B$.

For the second assertion, let $I$ be the face poset of the boundary of
a triangle and use the constant simplicial-set diagram with value a
point. Its strict colimit is a point, while its homotopy colimit is
$BI\cong S^1$. The category has six objects, the six identities and six
vertex--edge incidence arrows, and no composable pair of nonidentity
arrows. The normalized nerve chain calculation in
\texttt{code/AB\_checks.py} gives $H_0=\mathbb Z$, $H_1=\mathbb Z$,
and no higher homology. This explicitly prevents the proposed replacement.
\end{proof}

\begin{statement}{A.11}{Proved}{What restriction of an abstract derivator supplies}
For a derivator $\D$ in presheaf variance, $\D(P)$, its restriction
functors, both Kan extensions, and their pointwise comma formulas give
a well-defined theory of cellular coefficients on posets. The locally
constant subderivator has both adjoints as in A.9. If
$\D=\D_{\cC}$ and the geometric coefficient hypotheses of A.5 hold,
this theory is the homotopy category, at every diagram shape, of
constructible $\cC$-valued sheaves. Without an enhancement or a separately
defined sheaf theory, the phrase ``$\D$-valued topological sheaf'' is not
an additional consequence of Der~1--4.
\end{statement}
\begin{proof}
The first assertions are Der~3--4 and
\citet[Theorem 3.13]{Cisinski2009}. In the enhanced case, apply the
exodromy equivalence before taking homotopy categories and with an
additional indexing shape; the resulting restrictions and adjoints agree
by their universal properties. The final assertion concerns the data
available: Der~1--4 specify categories of coherent indexed diagrams and
Kan extensions, but do not specify a topology-valued sheaf construction
or mapping spaces. Defining cellular $\D$-coefficients by $\D(P)$ is
consistent and useful, but cannot prove an equivalence with an unspecified
category of topological sheaves.
\end{proof}

\subsection{The unrestricted raw exit calculation}

\begin{statement}{A.12}{Proved}{Raw exit paths of an arbitrary poset realization}
For every small poset $P$, the minimum-vertex map
$s_P:|NP|\to P_{\downarrow}$ is continuous. Its raw exit simplicial set,
formed using the upward-open target $P^{\op}$, is an $\infty$-category,
and the projection recording the stratum labels is a categorical
equivalence
\[
 \operatorname{Sing}^{P^{\op}}(|NP|,s_P)
       \longrightarrow N(P^{\op}).
\tag{A.12}
\]
Here $|NP|$ has its weak CW topology. This statement does not identify
ordinary constructible sheaves on that unrestricted realization.
\end{statement}
\begin{proof}
Each barycentric coordinate is continuous in the weak CW topology,
because its restriction to every closed simplex is continuous. The
inverse image of a downset is the union of the open vertex stars at its
elements, which proves continuity.

We recall the compactness property of a CW complex in the form needed
here. A compact subset meets only finitely many open cells. Indeed,
otherwise choose one point in each of countably many distinct open
cells. Every closed cell contains points from only finitely many of
these open cells, by closure finiteness. Consequently every subset of
the chosen point set has closed intersection with every closed cell,
and is closed in the CW complex by its weak topology. The point set
would therefore be an infinite closed discrete subset of a compact
space, a contradiction. Taking the closures of the finitely many cells
met by a compact subset places that subset in a finite subcomplex.
This is also the finite-subcomplex reduction used in
\citet[Corollary A.6.9]{LurieHA}.

Let $P_0$ run through the finite subsets of $P$, with their induced
orders. A finite subcomplex of $|NP|$ is contained in $|NP_0|$ for some
such $P_0$: take all vertices of its finitely many cells. Every singular
exit simplex has compact image and hence factors through $|NP_0|$.
Its stratum labels also belong to $P_0$, since every vertex of the
domain simplex has the prescribed label of its image. Thus there is
an equality of simplicial sets
\[
 \operatorname{Sing}^{P^{\op}}(|NP|,s_P)
 =\colim_{P_0\subset P,\;|P_0|<\infty}
       \operatorname{Sing}^{P_0^{\op}}(|NP_0|,s_{P_0}).
\tag{A.12a}
\]
The transition maps are monomorphisms. Any inner horn in this filtered
union occurs at a finite stage, where A.4 supplies a filler. Hence the
union is an $\infty$-category.

For completeness, the same finite-data observation verifies that the
union of the projections in A.4 is a categorical equivalence. Any two
objects occur together at a finite stage. The right mapping simplicial
sets between them are filtered unions of the corresponding right
mapping simplicial sets at these finite stages: their $n$-simplices
and defining endpoint conditions use only the finite simplex
$\Delta^{n+1}$. These are Kan complexes, and their filtered unions
compute their filtered colimits as spaces, since each horn, sphere,
and homotopy uses only finitely many simplices. A.4 identifies each
stage with the appropriate mapping space in $N(P_0^{\op})$. Thus the
resulting mapping space is contractible if the two stratum labels
$p,q$ satisfy $q\leq p$, and is empty otherwise. Every label $p$ occurs
at the vertex $p$, so the projection is also essentially surjective.
This proves (A.12), including mapping spaces. The argument concerns
finite domains of singular simplices; it does not imply that a
constructible sheaf is recovered from these simplices without an
exodromy theorem.
\end{proof}
